% Generated draft: requires mathematical and source review.
\documentclass[11pt]{article}
\usepackage[margin=25mm]{geometry}
\usepackage{fontspec}
\setmainfont{Latin Modern Roman}
\usepackage{amsmath,amssymb,mathtools,unicode-math}
\setmathfont{Latin Modern Math}
\usepackage{enumitem,needspace,xurl,hyperref}
\setlength{\parindent}{0pt}
\setlength{\parskip}{5pt}
\setlength{\emergencystretch}{3em}
\newcommand{\N}{\mathbb{N}}
\newcommand{\Z}{\mathbb{Z}}
\newcommand{\Q}{\mathbb{Q}}
\newcommand{\R}{\mathbb{R}}
\newcommand{\C}{\mathbb{C}}
\newcommand{\F}{\mathbb{F}}
\newcommand{\A}{\mathbb{A}}
\newcommand{\PP}{\mathbb{P}}
\newcommand{\E}{\mathbb{E}}
\newcommand{\eps}{\varepsilon}
\newcommand{\dd}{\,\mathrm{d}}
\begin{document}
\section*{UnsolvedMath ID: 700806}
\textbf{Source problem code:} catalogue entry
\subsection{Definitions and mathematical statement}
A group \(G\) is \emph{locally graded} if every nontrivial finitely generated subgroup \(K\leq G\) has a proper subgroup of finite index. It has \emph{finite special rank} if there exists an integer \(r\geq 0\) such that every finitely generated subgroup of \(G\) can be generated by at most \(r\) elements.

A group \(H\) is \emph{hyperabelian} if there exist an ordinal \(\lambda\) and an ascending series \((H_\alpha)_{\alpha\leq\lambda}\) of normal subgroups of \(H\) such that
\[
H_0=\{1\},\qquad H_\lambda=H,\qquad
H_{\alpha+1}/H_\alpha\text{ is abelian}\quad(\alpha<\lambda),
\]
and \(H_\beta=\bigcup_{\alpha<\beta}H_\alpha\) for every nonzero limit ordinal \(\beta\leq\lambda\). A group is \emph{almost hyperabelian} if it contains a hyperabelian subgroup of finite index.

\textbf{Question (N. S. Chernikov, 1990).} Is it true that for every locally graded group \(G\) of finite special rank there exists a subgroup \(H\leq G\) such that
\[
[G:H]<\infty\qquad\text{and}\qquad H\text{ is hyperabelian}?
\]

The supplied source states that the answer is affirmative when \(G\) is additionally \emph{periodic}, meaning that every element of \(G\) has finite order. It cites \emph{Ukrain. Math. J.}, volume 42, no. 7 (1990), pages 855--861. This special-case assertion is retained as supplied background; its reference is incomplete and has not been independently verified.
\subsection{Short English statement}
Does every locally graded group whose finitely generated subgroups have a uniform finite bound on their numbers of generators contain a hyperabelian subgroup of finite index? The source attributes the question to N. S. Chernikov and reports an affirmative answer for periodic groups.
\subsection{Sources}
\begin{itemize}
\item Local UnsolvedMath registry, permanent ID \(700806\), exact source code \texttt{catalogue entry}, title \emph{Kourovka Notebook Problem 11.114}; supplied upstream revision \texttt{d371ee917ac9647fbb06cd6fc529deec6f3dda82}.
\item \emph{The Kourovka Notebook (Unsolved Problems in Group Theory)}, original problem 11.114, 11th issue (1990), proposed by N. S. Chernikov. The supplied record identifies the 21st-issue edition (2026), \texttt{arXiv:1401.0300v48}, and names E. I. Khukhro and V. D. Mazurov as editors of the collection.
\item For the reported periodic case: \emph{Ukrain. Math. J.}, \textbf{42}, no. 7 (1990), 855--861. The supplied record gives no article title or authors for this citation.
\item The registry labels the problem \texttt{open}. This is a supplied status label, not an independently verified assessment of the literature.
\end{itemize}
\end{document}
